% Korean edition of sections/s0_cover.tex.
% Cover: same coordinate grid as the 7709 report cover.tex (mm from page top-left, text baselines).
% Title Arial Bold 14pt, subtitle 10pt grey; right column Quant Team + authors (name 8pt bold, e-mail +3.9mm, pitch 9.5mm; no phone numbers in the public version).
% \covertitlelines (1 or 2) must match the typeset title; the subtitle moves down 6.88mm per extra title line.
\thispagestyle{dartpage}
\begin{tikzpicture}[remember picture,overlay,x=1mm,y=-1mm,shift={(current page.north west)}]
  \node[anchor=base west,inner sep=0,text width=136.2mm,align=left] at (14.54,32.69)
    {\fontsize{20}{24}\selectfont\bfseries\addfontfeatures{LetterSpace=0.64}\color{navy}\RaggedRight\hyphenpenalty=10000 \exhyphenpenalty=10000 \relax\covertitle};
  \node[anchor=base west,inner sep=0,text width=136.2mm,align=left] at (14.54,{41.6+6.88*(\covertitlelines-1)})
    {\fontsize{10}{12}\selectfont\addfontfeatures{LetterSpace=0.41}\color{subgrey}\coversubtitle};
  \node[anchor=base west,inner sep=0] at (154.57,31.55) {\kofont\bfseries\fontsize{10}{10}\selectfont\addfontfeatures{LetterSpace=4.2}\color{navy}Quant Team};
  \fill[navy] (154.13,33.15) rectangle (197.73,33.49);
  \def\analystlist{38.02/Choi Munseok/chlanstjr0301@gmail.com,
                   47.52/Choi Yeayun/dbs6463@naver.com}
  \foreach \y/\n/\e in \analystlist {
    \node[anchor=base west,inner sep=0] at (154.57,\y) {\kofont\bfseries\fontsize{8}{8}\selectfont\addfontfeatures{LetterSpace=3.0}\color{black}\n};
    \node[anchor=base west,inner sep=0] at (154.57,\y+3.9) {\kofont\fontsize{7}{7}\selectfont\addfontfeatures{LetterSpace=2.1}\color{body}\e};
  }
  \node[anchor=base west,inner sep=0] at (14.54,280.6) {\fontsize{7.5}{9}\selectfont\color{subgrey}코드: \href{https://github.com/chlanstjr0301/pytimefolio}{\code{github.com/chlanstjr0301/pytimefolio}}};
\end{tikzpicture}%
\vspace*{\dimexpr 6.88mm*(\covertitlelines-1)+26.4mm\relax}%
\par
\newcommand{\coverblock}[2]{{\fontsize{9.5}{11.8}\selectfont\bfseries\color{navy}\RaggedRight #1\par}\vspace{0.5pt}{\fontsize{9}{11.4}\selectfont #2\par}\vspace{2.5pt}}
\coverblock{이 설계 문서는 래퍼의 모든 결정을 측정값이나 명시된 규칙에 연결한다}{이 문서는 \code{pytimefolio}의 소프트웨어 설계 문서다. 대회 서버가 실제로 어떻게 동작하는지(\ref{sec:layers}장), 거기서 나온 설계(\ref{sec:design}장), 실서버 테스트(\ref{sec:live}장)를 기록한다. \figref{summary}는 각 계층과 그에 따른 결정을 정리한다. 타임폴리오는 Road to Fund Manager 대회의 API를 공개하지 않으므로, 래퍼는 대회 웹 클라이언트가 보내는 요청을 그대로 보낸다. 우리는 이 요청을 재현하고 2026년 11월 24일 서버의 응답을 측정했다. 주최측은 이메일로 자동 주문을 금지하지 않으며, 하루 100건 이상 주문하는 IP 주소는 차단될 수 있다고 밝혔다 [5]. 그래서 래퍼는 하루 80건에서 멈춘다.}
\coverblock{측정값이 설계를 정한다: 연결 하나를 재사용하고, 토큰 하나를 쓰며, 3\,s마다 폴링한다}{\ref{sec:layers}장은 프로토콜 5개 계층을 측정한다. REST 요청은 재사용한 HTTP/2 연결에서 16.0\ms{}, 새 연결에서 106.5\ms{}가 걸린다. 액세스 토큰은 900\,s 동안 유효하다. 로그인 때 받은 쿠키 묶음은 7일간 유효하며 갱신이나 로그아웃으로 바뀌지 않으므로, 래퍼는 쿠키 묶음을 자격 증명으로 보호한다. 실시간 피드는 요청하지 않으면 아무것도 보내지 않으므로 래퍼가 3\,s마다 요청한다. 직전 응답의 스탬프(\code{t})를 정확히 되돌려 보낸 요청은 바뀐 필드만 담은 델타를 받는다(세션 중 1,053\textasciitilde 1,912바이트, 스냅샷은 11,344바이트). 그 밖의 요청은 전체 스냅샷을 받는다.}
\coverblock{주문 ID가 돌아오지 않으므로, 서버가 받았을 수 있는 주문은 절대 다시 보내지 않는다}{주문(\ref{sec:orders}절)은 REST로 넣고 피드에서 확인하며, \ref{sec:design}장은 이 안전 규칙을 중심으로 래퍼를 설계한다. 래퍼가 주문을 한 번 더 보내는 경우는 둘뿐이다. 연결이 아예 맺어지지 않았을 때와 401을 받았을 때다. 401은 ASP.NET이 주문 처리기를 실행하기 전에 돌려준다. 그 밖에는 15\,s 동안 피드에서 주문을 찾고, 다음으로 주문 목록에서 찾으며, 그래도 없으면 예외를 던져 호출자가 웹 페이지에서 확인하게 한다. 주문 전마다 피드의 나이(26\,s 미만), 영업일, 종목 한도를 검사한다. \code{plan()}은 목표 비중의 섹터 한도와 합계 한도를 검사한다.}
\coverblock{실서버 테스트에서 결함 2개가 드러나 고쳤고, 주식 수는 내림된다는 것을 확인했다}{\ref{sec:live}장의 실서버 테스트에서 시작, 피드, 매수, 매도가 모두 동작했다. 실서버에서 일부러 일으킬 수 없는 실패는 오프라인 테스트 31개가 다룬다. 실서버 테스트에서 결함 2개가 드러났다. \code{plan()}이 합계 139\%인 목표를 받아들였고, \code{guidelines()}가 잘못된 키로 섹터 한도를 찾았다. 1.76주에 해당하는 0.04\% 주문을 두 번 보냈는데, 서버는 두 번 모두 주문 주식 수(\code{qty2ord})를 1로 정했다. 반올림이 아니라는 뜻이므로 \code{plan()}은 1주 미만 주문을 뺀다.}
\vspace{2pt}

\figinlinetrue% Korean edition of figs/summary.tex.
\fig[wide]{summary}{계층마다 실서버에서 측정했고, 각 결과가 부록 B의 결정을 정한다}[절별 주요 결과와 그것이 뒷받침하는 결정(부록~\ref{app:decisions}의 D 번호)]{%
{\tabfont\begin{tabularx}{\linewidth}{>{\RaggedRight\arraybackslash}p{27mm} >{\RaggedRight\arraybackslash}X >{\RaggedRight\arraybackslash}p{29mm}}
\thd{계층} & \thd{결과} & \thd{결정}\\
\textbf{\ref{sec:transport} 전송} & 요청은 새 HTTP/2 연결에서 106.5\ms{}, 재사용 연결에서 16.0\ms{}가 걸린다. 서버는 15\textasciitilde 30\,s 유휴 후 연결을 닫는다 & D1, D2, D18\\\rowsep
\textbf{\ref{sec:auth} 로그인 세션} & 액세스 토큰은 900\,s 유효하며 만료 후 300\,s까지 받아준다. 쿠키 묶음은 7일 유효하고, 갱신으로 연장되지 않으며, 저장한 사본은 \code{Logout} 뒤에도 계속 갱신된다 & D3\textasciitilde D5, D16, D17, D26, D35, D36\\\rowsep
\textbf{\ref{sec:queries} 조회} & 모든 조회는 Bearer 헤더로 인증된다. 거래일이 아닌 날짜는 HTTP~200과 함께 다른 날의 데이터를 돌려준다 & D6, D7, D27\\\rowsep
\textbf{\ref{sec:feed} 실시간 피드} & 피드는 요청할 때만(3\,s마다) 답한다. \code{t}를 정확히 되돌리면 11,344바이트 스냅샷 대신 1,053\textasciitilde 1,912바이트 델타를 받는다. 열린 WebSocket은 토큰보다 오래 산다 & D8, D9, D14, D30, D31\\\rowsep
\textbf{\ref{sec:orders} 주문} & \code{AddOrder}는 주문 ID를 돌려주지 않는다. \code{wei}는 주문 크기다(0.3\%가 7주를 요구). 끝난 주문은 \code{Done}이 된다. 7주가 7분 37초에 체결됐다 & D11, D13, D15, D29, D32\\\rowsep
\textbf{\ref{sec:design} 래퍼 설계} & asyncio로 HTTP를 직접 보내 로그인 세션 하나를 유지한다. 전략과 시장 데이터는 호출자 몫이다. 하루 100건 이상 주문하는 IP 주소는 차단될 수 있다 & D10, D12, D34\\\rowsep
\textbf{\ref{sec:live} 실서버 테스트} & 시작, 피드, 매수, 매도가 동작했다. \code{qty2ord}는 내림한다. 결함 2개를 고쳤다. 15\,s 주문 검색은 느린 서버에서 시험하지 않았고, 한도는 목표에 대해서만 검사한다 & D19\textasciitilde D25, D28, D33\\
\end{tabularx}}}{2026년 11월 24일 05:38\textasciitilde 10:34 KST, 계정 하나로 측정. \ref{sec:transport}절: curl은 09:00 개장 전, Python은 장중. 3\,s 폴링은 10월 28일 기록에서.}{DART Quant Team}
\figinlinefalse

\clearpage
